\documentclass[11pt]{article}

\usepackage[margin=1in]{geometry}
\usepackage{amsmath}
\usepackage{booktabs}
\usepackage{array}
\usepackage{graphicx}
\usepackage{hyperref}

\title{From Trace Comparison to Parameter Optimization:\\The Machine-Learning Plan}
\author{ML-assisted optimization of CCD operating parameters --- working note}
\date{October 2026}

\newcommand{\code}[1]{\texttt{#1}}

\begin{document}
\maketitle

\begin{abstract}
The trace-comparison tools (companion note \code{trace\_fit\_tools.tex}) tell whether two raw video traces differ.
The project goal is larger: use features of a short trace to find operating parameters for which a calibrated
image has low readout noise $\sigma_e$, low dark current $\lambda$ and a known gain $g$. This note describes how that
would be trained and used, what data it needs, and what is still unproven. Nothing described here has been trained yet.
The figures are schematic diagrams of the method, not results.
\end{abstract}

\section{Goal and status}

\paragraph{Goal.} For a skipper CCD, find the settings (clock timing and levels, biases, number of skips, \ldots) that give the
best calibrated image: small $\sigma_e$ and small $\lambda$, with a known gain $g$. The trace is a cheap measurement
(about one second); the calibrated image is the expensive ground truth (minutes to hours).

\paragraph{What exists.}
\begin{itemize}
  \item Scripts that take a packet (one trace per channel plus one image, at identical settings) and scan many settings
        with safety limits (\code{take\_packet.py}, \code{scan\_packets.py}).
  \item A decoder, region splitter and 39 comparison tools (\code{decode\_trace.py}, \code{trace\_regions.py}, \code{compare\_trace.py}).
  \item A test on real data: the tools separate two amplifiers read with identical parameters, and repeats of one amplifier are
        distinguished from a different amplifier only when judged against the run-to-run spread.
\end{itemize}

\paragraph{What does not exist yet.} A scan with different operating parameters, labels ($\sigma_e$, $g$, $\lambda$) from calibrated
images for those settings, and any trained model. All traces analysed so far have identical operating parameters.

\section{What is being optimized}

The quantities to optimize are the \emph{operating parameters} themselves, in three groups:
\begin{center}
\begin{tabular}{p{0.2\textwidth}p{0.45\textwidth}p{0.25\textwidth}}
\toprule
Group & Parameters & In the scan code now \\
\midrule
Bias voltages & \code{VSUB}, \code{VDD}, \code{VDRAIN}, \code{VR} (per amplifier) & yes \\
Rail voltages (clock levels) & high and low levels of the vertical clocks, horizontal clocks and transfer gates (\code{V*\_H/L}, \code{H*\_H/L}, \code{TG\_*\_H/L}) & yes; the summing-well, output-gate, reset-gate and dump-gate rails are not yet \\
Timing & transfer hold times (\code{delay\_Va/Vb/Ha/Hb/Hc}); the integration, reset, summing-well and output-gate delays & transfer delays yes; the other delays are not yet \\
\bottomrule
\end{tabular}
\end{center}
The number of skips and the binning are held fixed for now (500 skips, one binning), so that images stay comparable.

\paragraph{Consequences for the method.}
\begin{itemize}
  \item The settings are the \emph{inputs of the surrogate model} and the \emph{search space of the optimizer}: the model learns
        settings (plus trace features) $\rightarrow$ $\sigma_e$, $g$, $\lambda$, and the optimizer chooses the next settings. Recovering which
        settings changed from a trace is a secondary, diagnostic task (fault detection, annotating old data).
  \item The space is high-dimensional (tens of parameters, because every clock pin has a high and a low level) while the number of
        labeled packets is a few hundred. It has to be reduced: screening scans to drop unimportant parameters, and a
        reparametrization such as one reference voltage plus offsets, using the group's empirical relations. Your proposal expects an effective
        dimension of 8 to 12.
  \item The parameters are not independent: the CCD works only if certain bias and clock relationships hold. They must be encoded as
        constraints (\code{relational\_constraints} in \code{ccd\_limits.py} is still an empty placeholder) so that the optimizer never proposes a forbidden
        combination.
  \item Changes are made from the best configuration so far in bounded steps (voltages are ramped in steps no larger than \code{max\_step}),
        so the optimizer works in a trust region around the incumbent, not in jumps.
  \item Every packet must store the settings that were applied, the settings of the reference it is compared with, the change from that
        reference, the previous settings and the time since the previous acquisition, and the values the board reports back. At present
        only the commanded settings are stored in the manifest; the other fields are still to be added.
\end{itemize}

\section{Training data: the labeled packet}

To learn ``trace features $\rightarrow$ $\sigma_e$, $g$, $\lambda$'' the model needs, for every operating point, \emph{both} a trace
and a calibrated image taken at the same settings. The settings are known (we set them), the tool values come from the trace,
and the physics values come from the image through the WADERs calibration (Figure~\ref{fig:example}).

\begin{figure}[htbp]
  \centering
  \includegraphics[width=0.85\textwidth]{../plots/ml_example.png}
  \caption{One labeled training example. The input is the settings plus the tool values of the trace; the answer is the read noise,
  gain and dark rate measured from the calibrated image taken at the same settings.}
  \label{fig:example}
\end{figure}

Only parts of the problem need image labels:
\begin{center}
\begin{tabular}{p{0.45\textwidth}cp{0.3\textwidth}}
\toprule
Quantity to predict & Image label needed? & Why \\
\midrule
Diagnostic only: does a trace match the reference, which settings changed & no & the settings are known \\
Read noise $\sigma_e$, gain $g$, dark current $\lambda$ & \textbf{yes} & measured only from calibrated images \\
Unlabeled structure (anomaly detection, self-supervised features) & no & any trace can be used \\
\bottomrule
\end{tabular}
\end{center}

\subsection{What each label costs}
\begin{itemize}
  \item \textbf{$\sigma_e$} is cheap: about 1250 empty pixels (the overscan) give 2\%. The image must use the same number of skips as
        the images the optimization is for.
  \item \textbf{$g$} needs pixels of known charge. At low dark current the 1\,e$^-$ peak has few events, so a short image with a few
        seconds of exposure may not give a gain. The gain depends on the voltages but not on the number of skips, so it can be measured
        once per voltage set with a dedicated higher-charge image and reused.
  \item \textbf{$\lambda$} is expensive: about $10^6$ pixels for a 10\% measurement at $10^{-4}\,e^-$/pixel. It should be budgeted for a
        subset of the points (order 100 to 300), at two exposure times to separate exposure-dependent dark current from
        exposure-independent spurious charge.
\end{itemize}

\subsection{Conditions on the pairing}
\begin{itemize}
  \item Same settings, same amplifier, taken close in time: the CCD state changed between consecutive runs in the data so far.
  \item One image gives four labels (one per amplifier extension) and the scan takes one trace per channel; the channel-to-extension
        mapping must be kept explicit.
  \item The WADERs output does not store the fit $\chi^2$, quality flags or software version; the calibration wrapper must save
        them with each packet, together with the configuration file.
  \item The label noise floor must be measured: several repeats of the reference, each with an image. No model can predict better
        than this scatter.
\end{itemize}

\section{Training}

Figure~\ref{fig:training} shows the loop: the model starts with arbitrary settings of its adjustable parameters, predicts the physics values for each packet,
and is corrected by the difference to the image-derived values.

\begin{figure}[htbp]
  \centering
  \includegraphics[width=0.9\textwidth]{../plots/ml_training.png}
  \caption{The training loop. The model predicts the physics values from the input, the error against the image-derived values is used
  to adjust it, and the process is repeated. Held-out settings and amplifiers, never used to adjust the model, give the score.}
  \label{fig:training}
\end{figure}

\paragraph{Models, from simple to ambitious.}
\begin{center}
\begin{tabular}{p{0.28\textwidth}p{0.42\textwidth}p{0.2\textwidth}}
\toprule
Model & Use & Data needed \\
\midrule
Reference-consistency test (robust Mahalanobis distance) & is a trace consistent with the reference repeats & available now \\
Classifier two-sample test (logistic regression, trees; grouped cross-validation) & can a model tell trace sets apart on unseen runs & available now \\
Ridge / partial least squares / gradient-boosted trees & settings and tool values $\rightarrow$ $\sigma_e$, $g$, $\lambda$; the directions in which parameters move the tools & scan with labels \\
Gaussian process surrogate & settings $\rightarrow$ physics values with calibrated uncertainty, for optimization & scan with labels \\
Neural network on folded cycles & learn features directly & thousands of packets, many amplifiers \\
\bottomrule
\end{tabular}
\end{center}

\paragraph{Validation rules.}
\begin{itemize}
  \item Split by run, never by pixel or skip: pixels of one run share run-wide shifts, and a model validated on them would look good
        and mean nothing.
  \item Hold out whole settings values and, as the hardest test, whole amplifiers; report the error against how many packets of the
        new amplifier the model was allowed to see.
  \item Work with each amplifier's change relative to its own reference where possible, so the fixed amplifier signature
        does not have to be learned.
  \item Keep row-start pixels and charge-contaminated pixels out of the inputs (\code{usable\_pixels}).
  \item Cycle and block lengths are measured exactly and need no learning; the model works on length-independent features.
\end{itemize}

\section{Using the model to optimize}

\begin{figure}[htbp]
  \centering
  \includegraphics[width=0.9\textwidth]{../plots/ml_optimize.png}
  \caption{The optimization loop. A trace at candidate settings gives tool values; the trained model predicts $\sigma_e$, $g$ and $\lambda$ with
  an uncertainty; an optimizer proposes better settings. When the model is unsure, a calibrated image is taken and becomes a new
  training example.}
  \label{fig:optimize}
\end{figure}

Figure~\ref{fig:optimize} shows how the trained model is used. This is multi-fidelity optimization: many cheap trace measurements, a few expensive images. The optimizer is Bayesian optimization on top of
the surrogate model, with the safety limits (\code{limits.toml}) enforced outside the model and a human approval step early on.
A stopping rule: the targets are met with calibrated confidence, or the expected improvement falls below the reproducibility floor.

\subsection{The trade-off being searched}

Figure~\ref{fig:tradeoff} sketches what the optimizer is balancing.

\begin{figure}[htbp]
  \centering
  \includegraphics[width=0.75\textwidth]{../plots/ml_tradeoff.png}
  \caption{Schematic: more skips lower the read noise but lengthen the readout and raise dark current and spurious charge, so the
  best settings are a compromise. The optimization is multi-objective; the experiment picks a point on the resulting front according to its science goal.
  The curves are illustrative, not measured.}
  \label{fig:tradeoff}
\end{figure}

\section{What the first scan should provide}
\begin{itemize}
  \item The reference configuration at least five times, interleaved through the scan, each with an image, to measure run-to-run and
        label scatter.
  \item Several values of each parameter to be recovered, not a single step, plus some combinations of parameters changed together:
        two parameters can move the same tool and one-at-a-time changes cannot separate them.
  \item All four amplifiers, ideally more than one module.
  \item Images for $\sigma_e$ at the target number of skips, a gain image per voltage set, and long exposures for $\lambda$ at a subset of points.
  \item A pause or flush between acquisitions, and a record of the time since the last acquisition and the previous settings.
\end{itemize}

\section{Risks and open questions}
\begin{itemize}
  \item The tools react to the amplifier much more than to the run-to-run noise; whether parameter changes move them by more than about 0.5\%
        of the edge heights (about three times the run-to-run spread) is unknown.
  \item Few runs and few pixels limit every statistic; the smallest reportable $p$-value with 8 usable pixels is $1.6\times10^{-4}$.
  \item The CDS level changed between back-to-back acquisitions; the cause (real charge in the pixels) is not confirmed.
  \item The first-skip transient and the anti-correlated skip noise mean simple white-noise models do not describe the noise.
  \item Generalizing to amplifiers never seen is the hardest goal; a one-reference-packet scheme or a few-shot adaptation is the realistic path.
  \item The relationships between biases and clock levels required for proper CCD operation are not yet encoded in \code{ccd\_limits.py}.
\end{itemize}

\section{Where this lives in the code}
\begin{center}
\begin{tabular}{ll}
\toprule
Item & File \\
\midrule
Packet and scan acquisition      & \code{src/take\_packet.py}, \code{src/scan\_packets.py}, \code{src/ccd\_limits.py} \\
Decoding, regions, tools         & \code{src/decode\_trace.py}, \code{src/trace\_regions.py}, \code{src/compare\_trace.py} \\
Walkthrough notebook             & \code{notebooks/trace\_comparison\_walkthrough.ipynb} \\
These diagrams                   & \code{src/plot\_ml\_diagrams.py} \\
Packet processing (decode, WADERs, grouping) & designed, not yet written \\
Models                           & not yet written \\
\bottomrule
\end{tabular}
\end{center}

\end{document}
